% GENERATED FILE. Edit canonical lessons, then run npm run generate:books.
% canonical-source-hash: fnv1a64:a750ac606dd3fbff
% canonical-lessons: ZH-C317-qingnian, ZH-C317-chengnian, ZH-C317-hunyin, ZH-C317-jiating, ZH-C317-shixian, ZH-R317-review-313-317, ZH-R317-review-final-317

\chapter{Looking Back}
\label{ch:zh-looking-back}
\hlchaptermodality{317}

\begin{chapteropening}
\textbf{By the end of this chapter:} \emph{I can talk about growing up, marriage and family, and doing things ahead of time.}

Complete the last lesson of chapter 317: I can talk about growing up, marriage and family, and doing things ahead of time.
\end{chapteropening}

\section[qīngnián]{qīngnián — young people; youth}

\label{lesson:ZH-C317-qingnian}



\begin{quote}
Before the new one: say the Mandarin for life from birth to death, then the Mandarin for as you would naturally expect.
\end{quote}

\subsection*{You'll want to know: qīngnián}
\emph{qīngnián} — \textquotedblleft{}young people; youth\textquotedblright{}.

\begin{quote}
\emph{tā shì yí ge hěn yǒu nénglì de qīngnián} — He's a very capable young man.
\end{quote}

Its characters come later; for now it is a word for the ear.

\begin{cousinweb}
\emph{qīngnián} is built from \emph{qīng} 'green' and \zh{年} 'years', both inherited from Middle Chinese, \emph{tsheng} and \emph{nen} (Guangyun, in Baxter's transcription), which glosses the first as the colour of the east: green years.
\end{cousinweb}

\begin{grammarlens}[title={What you've built}]
One more word for telling what happened.
\end{grammarlens}

\subsection*{Your turn}
\begin{itemize}
\raggedright
  \item \emph{Say it:} \emph{qīngnián}
  \item \emph{Say it:} \emph{qīngnián}, once more
  \item \emph{Say it:} \emph{tā shì yí ge hěn yǒu nénglì de qīngnián} — He's a very capable young man.
  \item \emph{Recall:} say the Mandarin for life from birth to death, then the Mandarin for as you would naturally expect, then say \emph{qīngnián} again
\end{itemize}

\begin{tcolorbox}[breakable,colback=teal!4,colframe=teal!35!black,title={Before you move on}]
What does qīngnián mean? (\textquotedblleft{}Young people; youth\textquotedblright{}.) Say it once more.
\end{tcolorbox}
\section[chéngnián]{chéngnián — to come of age; fully grown}

\label{lesson:ZH-C317-chengnian}



\begin{quote}
Before the new one: say the Mandarin for as you would naturally expect, then the Mandarin for young people.
\end{quote}

\subsection*{You'll want to know: chéngnián}
\emph{chéngnián} — \textquotedblleft{}to come of age; fully grown\textquotedblright{}.

\begin{quote}
\emph{háizi chéngnián yǐhòu, jiù líkāi le jiā} — Once the children were grown, they left home.
\end{quote}

Its characters come later; for now it is a word for the ear.

\begin{cousinweb}
\emph{chéngnián} is built from \emph{chéng} 'to complete' and \zh{年} 'years', both inherited from Middle Chinese, \emph{dzyeng} and \emph{nen} (Guangyun, in Baxter's transcription), which glosses the first 'to finish' and 'to make'.
\end{cousinweb}

\begin{grammarlens}[title={What you've built}]
One more word for telling what happened.
\end{grammarlens}

\subsection*{Your turn}
\begin{itemize}
\raggedright
  \item \emph{Say it:} \emph{chéngnián}
  \item \emph{Say it:} \emph{chéngnián}, once more
  \item \emph{Say it:} \emph{háizi chéngnián yǐhòu, jiù líkāi le jiā} — Once the children were grown, they left home.
  \item \emph{Recall:} say the Mandarin for as you would naturally expect, then the Mandarin for young people, then say \emph{chéngnián} again
\end{itemize}

\begin{tcolorbox}[breakable,colback=teal!4,colframe=teal!35!black,title={Before you move on}]
What does chéngnián mean? (\textquotedblleft{}To come of age; fully grown\textquotedblright{}.) Say it once more.
\end{tcolorbox}
\section[hūnyīn]{hūnyīn — marriage; the bond of marriage}

\label{lesson:ZH-C317-hunyin}



\begin{quote}
Before the new one: say the Mandarin for young people, then the Mandarin for to come of age.
\end{quote}

\subsection*{You'll want to know: hūnyīn}
\emph{hūnyīn} — \textquotedblleft{}marriage; the bond of marriage\textquotedblright{}.

\begin{quote}
\emph{tāmen de hūnyīn yǒu wèntí} — Their marriage is in trouble.
\end{quote}

Its characters come later; for now it is a word for the ear.

\begin{cousinweb}
\emph{hūnyīn} is built from \emph{hūn} and \emph{yīn}, both inherited from Middle Chinese, \emph{xwon} and \emph{'jin} (Guangyun, in Baxter's transcription), which glosses each with the pair itself, 'marriage', and explains \emph{hūn} by the evening hour of old weddings.
\end{cousinweb}

\begin{grammarlens}[title={What you've built}]
One more word for telling what happened.
\end{grammarlens}

\subsection*{Your turn}
\begin{itemize}
\raggedright
  \item \emph{Say it:} \emph{hūnyīn}
  \item \emph{Say it:} \emph{hūnyīn}, once more
  \item \emph{Say it:} \emph{tāmen de hūnyīn yǒu wèntí} — Their marriage is in trouble.
  \item \emph{Recall:} say the Mandarin for young people, then the Mandarin for to come of age, then say \emph{hūnyīn} again
\end{itemize}

\begin{tcolorbox}[breakable,colback=teal!4,colframe=teal!35!black,title={Before you move on}]
What does hūnyīn mean? (\textquotedblleft{}Marriage; the bond of marriage\textquotedblright{}.) Say it once more.
\end{tcolorbox}
\section[jiātíng]{jiātíng — household; a family unit}

\label{lesson:ZH-C317-jiating}



\begin{quote}
Before the new one: say the Mandarin for to come of age, then the Mandarin for a marriage.
\end{quote}

\subsection*{You'll want to know: jiātíng}
\emph{jiātíng} — \textquotedblleft{}household; a family unit\textquotedblright{}.

\begin{quote}
\emph{wǒmen shì yí ge kuàilè de jiātíng} — We're a happy family.
\end{quote}

Its characters come later; for now it is a word for the ear.

\begin{cousinweb}
\emph{jiātíng} is built from \zh{家} 'home' and \emph{tíng} 'a courtyard', both inherited from Middle Chinese, \emph{kae} and \emph{deng} (Guangyun, in Baxter's transcription), which glosses the second 'a gate yard': the home and its yard.
\end{cousinweb}

\begin{grammarlens}[title={What you've built}]
One more word for telling what happened.
\end{grammarlens}

\subsection*{Your turn}
\begin{itemize}
\raggedright
  \item \emph{Say it:} \emph{jiātíng}
  \item \emph{Say it:} \emph{jiātíng}, once more
  \item \emph{Say it:} \emph{wǒmen shì yí ge kuàilè de jiātíng} — We're a happy family.
  \item \emph{Recall:} say the Mandarin for to come of age, then the Mandarin for a marriage, then say \emph{jiātíng} again
\end{itemize}

\begin{tcolorbox}[breakable,colback=teal!4,colframe=teal!35!black,title={Before you move on}]
What does jiātíng mean? (\textquotedblleft{}Household; a family unit\textquotedblright{}.) Say it once more.
\end{tcolorbox}
\section[\texorpdfstring{\zh{事先} (shìxiān)}{事先 (shìxiān)}]{\zh{事先} (shìxiān) — beforehand; in advance}

\label{lesson:ZH-C317-shixian}



\begin{quote}
Before the new one: say the Mandarin for a marriage, then the Mandarin for a household.
\end{quote}

\subsection*{You'll want to know: \zh{事先}}
\textbf{\zh{事先}} — \emph{shìxiān} — \textquotedblleft{}beforehand; in advance\textquotedblright{}.

\begin{quote}
\textbf{\zh{他又没有事先说我们都生气了}} \emph{tā yòu méiyǒu shìxiān shuō, wǒmen dōu shēngqì le} — Once again he didn't say so beforehand, and we were all cross.
\end{quote}

\begin{cousinweb}
\textbf{\zh{事先}} (\emph{shìxiān}) is built from \zh{事} 'a matter' and \zh{先} 'before', both inherited from Middle Chinese, \emph{dzriH} and \emph{sen} (Guangyun, in Baxter's transcription), which glosses the second 'first and after': before the matter.
\end{cousinweb}

\begin{grammarlens}[title={What you've built}]
One more word for telling what happened.
\end{grammarlens}

\subsection*{Your turn}
\begin{itemize}
\raggedright
  \item \emph{Say it:} \emph{shìxiān}
  \item \emph{Say it:} \emph{shìxiān}, once more
  \item \emph{Say it:} \emph{tā yòu méiyǒu shìxiān shuō, wǒmen dōu shēngqì le} — Once again he didn't say so beforehand, and we were all cross.
  \item \emph{Recall:} say the Mandarin for a marriage, then the Mandarin for a household, then say \emph{shìxiān} again
\end{itemize}

\begin{tcolorbox}[breakable,colback=teal!4,colframe=teal!35!black,title={Before you move on}]
What does \zh{事先} mean? (\textquotedblleft{}Beforehand; in advance\textquotedblright{}.) Say it once more.
\end{tcolorbox}
\section[(dialogue)]{Looking back: from Moving House to Looking Back}

\label{lesson:ZH-R317-review-313-317}



\begin{quote}
Say the Mandarin for to renovate a flat and not until then.
\end{quote}

\subsection*{Your turn}
\begin{itemize}
\raggedright
  \item \emph{Say it:} to renovate a flat, not until then, to make a place your home, to hate to leave behind, to adjust to new surroundings
  \item \emph{Say it:} a funeral, in someone's lifetime, to comfort someone who is sad, a cemetery, to shed tears
  \item \emph{Say it:} little by little, it happened again, at ease in a place, to my surprise, the same as before
  \item \emph{Say it:} to register officially, to be familiar with, of that place, life from birth to death, as you would naturally expect
  \item \emph{Say it:} young people, to come of age, a marriage, a household, in advance
  \item \emph{Read it:} \zh{才} \zh{又} — each one aloud, then a word that holds it
\end{itemize}

\begin{tcolorbox}[breakable,colback=teal!4,colframe=teal!35!black,title={Before you move on}]
To renovate a flat? (\emph{zhuāngxiū}.) Not until then? (\textbf{\zh{才}}, \emph{cái}.) To make a place your home? (\emph{dìngjū}.) To hate to leave behind? (\emph{shěbude}.) To adjust to new surroundings? (\emph{shìyìng}.) A funeral? (\emph{zànglǐ}.) In someone's lifetime? (\textbf{\zh{生前}}, \emph{shēngqián}.) To comfort someone who is sad? (\emph{ānwèi}.) A cemetery? (\emph{mùdì}.) To shed tears? (\emph{liúlèi}.) Little by little? (\emph{jiànjiàn}.) It happened again? (\textbf{\zh{又}}, \emph{yòu}.) At ease in a place? (\textbf{\zh{自在}}, \emph{zìzai}.) To my surprise? (\emph{jìngrán}.) The same as before? (\emph{réngrán}.) To register officially? (\emph{dēngjì}.) To be familiar with? (\emph{shúxī}.) Of that place? (\textbf{\zh{当地}}, \emph{dāngdì}.) Life from birth to death? (\textbf{\zh{人生}}, \emph{rénshēng}.) As you would naturally expect? (\textbf{\zh{自然}}, \emph{zìrán}.) Young people? (\emph{qīngnián}.) To come of age? (\emph{chéngnián}.) A marriage? (\emph{hūnyīn}.) A household? (\emph{jiātíng}.) In advance? (\textbf{\zh{事先}}, \emph{shìxiān}.)
\end{tcolorbox}
\section[(dialogue)]{Looking Back — once more}

\label{lesson:ZH-R317-review-final-317}



\begin{quote}
Say the Mandarin for young people and to come of age.
\end{quote}

\subsection*{Your turn}
\begin{itemize}
\raggedright
  \item \emph{Say it:} young people, to come of age, a marriage, a household, in advance
\end{itemize}

\begin{tcolorbox}[breakable,colback=teal!4,colframe=teal!35!black,title={Before you move on}]
Young people? (\emph{qīngnián}.) To come of age? (\emph{chéngnián}.) A marriage? (\emph{hūnyīn}.) A household? (\emph{jiātíng}.) In advance? (\textbf{\zh{事先}}, \emph{shìxiān}.)
\end{tcolorbox}
